Rozwój modeli sztucznej inteligencji pędzi bez opamiętania. W ciągu kilku ostatnich lat wszyscy byliśmy świadkami niesamowitych postępów w dziedzinach uczenia maszynowego, głębokiego uczenia i przetwarzania języka naturalnego. Coraz więcej osób przekonuje się do różnego rodzaju modeli SI i używa ich w coraz szerszych obszarach swojego życia. Niestety wraz z tym wszystkim dobre modele niespecjalnie stają się mniejsze, wręcz przeciwnie. Wiele z nich jest ogromnych rozmiarów przez ciągle zwiększające się ilości parametrów i nieustannie komplikujące się architektury, co utrudnia ich przechowywanie, przenoszenie oraz wczytywanie do pamięci urządzenia. Dla przykładu, najnowszy state-of-the-art model z rodziny Qwen zajmuje kilkaset gigabajtów \cite{qwen2025technical}. Wspomniany problem dotyczy nie tylko przestrzeni pamięci, czy to dyskowej, czy sprzętowej, ale też ma wpływ na wydajność, szybkość, a także na koszty uruchomienia modelu. W odpowiedzi na ten problem inżynierowie pracują nad różnymi metodami kompresji modeli, które pozwalają na zmniejszenie ich rozmiaru o różnym stopniu straty jakości. Jedną z takich metod jest kwantyzacja, czyli proces, w którym wartości parametrów modelu są reprezentowane za pomocą mniejszej liczby bitów, co pozwala na znaczne zmniejszenie rozmiaru modelu. Kwantyzacja może pomagać zmniejszyć zajmowane miejsce na dysku przez model, ale jest specyficzną formą kompresji, która niekoniecznie wymaga kroku dekompresji. Oznacza to, że modele skompresowane za pomocą niektórych kwantyzacji mogą być bezpośrednio używane bez potrzeby wcześniejszego dekompresowania, co jest dużą zaletą w kontekście wydajności i szybkości działania. \\

W kolejnych rozdziałach pracy przejdę przez ważne opisy, terminy i definicje, które pozwolą zrozumieć dokładniej, czym jest kwantyzacja, jakie dokładnie są jej rodzaje, korzyści oraz koszty. Następnie przedstawiony zostanie mój autorski przykład implementacji kwantyzacji wraz z wynikami testów na różnego rodzaju modelach w różnych konfiguracjach kompresji. Ostatnia część jest poświęcona publicznym rozwiązaniom i formatom, z~których korzystają najbardziej popularne modele i korporacje.